% ---------------------------------------------------------------------------
% dropout_layer.tex -- the dropout layer, rungs 1 and 2: equations, checks, libraries.
% \input from simulator_figures.tex, after editing_layer.tex.
% Code:  rung 1  src/17_dropout.py (generator + cache writer), src/18_validate_dropout.py,
%                src/21_f4_followup.py, src/19/20 (library);
%        rung 2  src/22_silencing.py (generator), src/23_validate_silencing.py, src/24/25 (library).
% Numbers: ../results/validate_dropout.json (the PASS, seed 2102026),
%          ../results/f4_followup.json, ../results/dropout_library_summary.json,
%          ../results/validate_silencing.json (the PASS, seed 3102026),
%          ../results/silencing_library_summary.json.
% Theory record: notes/sciphy_notes.md S5 (rung 1) and S6 (rung 2).
% No figures yet: neither rung has been illustrated.
% ---------------------------------------------------------------------------

% ===========================================================================
\section{The dropout layer, rung 1: technical dropout}\label{sec:drop1}
% ===========================================================================

The editing layer says what each cell's tapes \emph{read}. The dropout layer says which of those
tapes we actually \emph{see}. In Park's data a large share of cell $\times$ tape entries is missing,
and earlier analyses established that the missingness has two components: a technical one (a tape
not captured in a cell) and a heritable one (a tape silenced in a lineage and lost in all its
descendants). The simulator adds them in two steps, so that each one's cost can be measured
separately:

\begin{center}\small
\begin{tabular}{@{}lp{0.42\textwidth}p{0.36\textwidth}@{}}
\toprule
rung & what is simulated & what the step from the rung below measures \\
\midrule
0 & tree and editing, every tape observed (\S\ref{sec:edit}) & the pure homoplasy floor \\
1 & $+$ technical dropout with \emph{no} lineage structure (this section) & what ordinary technical
      dropout manufactures \\
2 & $+$ heritable silencing, complete Dollo loss (\S\ref{sec:drop2}) & what the silencing
      established on Park actually costs \\
\bottomrule
\end{tabular}
\end{center}

Rung~1 has a second job. \textbf{No statistic in this project has been run on data known to contain
no silencing}, and the heritable-silencing statistics
(\texttt{2026-08\_park-compatibility/src/73}, \texttt{74}, \texttt{77}) need exactly that as a negative
control. Rung~1 is that dataset, and it is written in the formats those scripts already read.

\subsection*{Notation}

Several letters here are also used elsewhere in this document. The table says which meaning applies
in the dropout sections.

\begin{center}\small
\begin{tabular}{@{}lp{0.55\textwidth}l@{}}
\toprule
symbol & meaning & units / range \\
\midrule
$c$, $z$ & a captured cell (a tip of the tree), and a tape, $z=1,\dots,k$. & --- \\
$D_{cz}$ & true depth of tape $z$ in cell $c$ at harvest, from the editing layer & $0$--$N$ \\
$\mu_z$ & $r_z\LamT$: tape $z$'s total expected edits over the experiment, before saturation & edits / tape \\
$X_{cz}$ & technical missingness: $1$ if tape $z$ is not captured in cell $c$ & $\{0,1\}$ \\
$\alpha_c$, $\beta_z$ & cell capture offset and tape recovery offset; larger means worse captured or
      worse recovered. & logits \\
$\pi_{cz}$ & $\sigma(\alpha_c+\beta_z)$, the probability that $X_{cz}=1$;
      $\sigma(x)=1/(1+e^{-x})$ & probability \\
$\pbar$, $\pbar_z$ & mean technical missingness over the panel's tapes; and for tape $z$, averaged
      over cells & probability \\
$s_\alpha$, $s_\beta$ & spread of $\alpha_c$ across cells; scatter of $\beta_z$ among tapes of one
      class & logits \\
$u_z$ & closed-locus class of tape $z$: $1$ = closed (poorly read \emph{and} slow) & $\{0,1\}$ \\
$\phi$ & fraction of tapes that are closed & $[0,1)$ \\
$\mu_\beta$, $\Delta_\beta$ & overall recovery level (solved, \S\ref{sec:pbar}); closed-tape read
      penalty & logits \\
$r_{\rm cl}$, $s_r$ & speed of a closed tape relative to an open one; lognormal speed scatter &
      dimensionless \\
$Y_{cz}$, $O_{cz}$ & observed: $Y=1$ if the tape is read with at least one edit; $O=1-Y$ & $\{0,1\}$ \\
$R_c$, $R_{\min}$ & tapes observed in cell $c$, $\sum_zY_{cz}$; the QC threshold on it. & tapes \\
$\psi(R_c)$ & probability that cell $c$ keeps its clone barcode (ClonalBC retention) & probability \\
$A$, $t_A$, $W_A$ & last common ancestor of two cells, its time, and $W(t_A)$, the share of
      editing spent before it & ---, $[0,1]$, $[0,1]$ \\
\bottomrule
\end{tabular}
\end{center}

\subsection{The model, in four parts}\label{sec:drop1-model}

\paragraph{The idea.} A tape can be missing from a cell's readout for reasons that have nothing to
do with the cell's ancestry: the cell was poorly captured, or the tape sits at an integration site
that is poorly read. Rung~1 models exactly those reasons and nothing else. Every missing entry is
then explained by a property of the cell and a property of the tape, and \emph{the tree plays no
part} --- that is the definition of the rung.

\paragraph{The order things are generated in.} Each piece uses something the one before it produced:
\begin{enumerate}
  \item \textbf{the panel} --- every tape's class $u_z$, recovery offset $\beta_z$ and speed $r_z$
        (\S\ref{sec:latent});
  \item \textbf{editing} --- the editing layer of \S\ref{sec:edit}, run with tape $z$ at speed $r_z$,
        so its expected edit count over the experiment is $\mu_z=r_z\LamT$; this produces the true
        depths $D_{cz}$;
  \item \textbf{the mask} --- each cell's capture offset $\alpha_c$, then $X_{cz}$ from $\alpha_c$ and
        $\beta_z$;
  \item \textbf{censoring} --- combine $X_{cz}$ with $D_{cz}$ into what is observed, $Y_{cz}$;
  \item \textbf{filters} --- drop or unlabel cells by how many tapes they show.
\end{enumerate}
The panel comes first because the speed $r_z$ is needed by editing, and the recovery offset
$\beta_z$ by the mask; both are properties of the same tape, which is what lets them be coupled.

\paragraph{1. The technical mask.} \emph{The logit scale first,} since the mask is written on it. For
a probability $p$, $\operatorname{logit}(p)=\ln\!\big(p/(1-p)\big)$, the log of the odds; its inverse
is the logistic function $\sigma(x)=1/(1+e^{-x})$. A logit of $0$ is a probability of $0.5$, $-0.85$
is $0.30$, and $+1.15$ is $0.76$. Adding $1$ to a logit multiplies the odds by $e\approx2.72$, whatever
the starting point. Working on this scale lets effects be \emph{added} while the probability stays
inside $(0,1)$.

Each entry is missing with a probability set by one number for the cell and one for the tape:
\begin{equation}
  \Prob(X_{cz}=1)=\pi_{cz}=\sigma(\alpha_c+\beta_z),
  \qquad \alpha_c=s_\alpha a_c,\quad a_c\sim\mathcal N(0,1)\ \text{iid over cells.}
  \label{eq:mask}
\end{equation}
Term by term:
\begin{itemize}
  \item $\beta_z$ is the \emph{tape's} recovery offset: how hard tape $z$ is to read, in any cell. It
        is set by the panel (\S\ref{sec:latent}).
  \item $\alpha_c$ is the \emph{cell's} capture offset: how badly cell $c$ was captured, on every tape.
        Larger means more of its tapes are missing.
  \item $\pi_{cz}$ is the resulting probability that tape $z$ is missing from cell $c$.
\end{itemize}
In odds the model reads $\pi_{cz}/(1-\pi_{cz})=e^{\alpha_c}\,e^{\beta_z}$: the cell and the tape each
contribute a \emph{factor}, and neither changes the other's. A poorly captured cell is worse on every
tape by the same odds factor, and a poorly read tape is worse in every cell by the same factor.

\paragraph{The cell scatter $s_\alpha$.} $\alpha_c$ is written as a product of two parts. $a_c$ is a
standard normal draw: how many standard deviations cell $c$ sits from an average cell. $s_\alpha$, in
logits, converts that into a capture offset, so $s_\alpha$ is the standard deviation of $\alpha_c$
across cells --- \emph{how unequal cells are in capture.} At $s_\alpha=0$ every cell is captured
alike and a tape's missingness is the same in all of them. At $s_\alpha=0.7$, the value in the
validation, a tape that is missing $30\%$ of the time in an average cell is missing $46\%$ of the time
in a cell one standard deviation worse ($a_c=+1$) and $18\%$ in one a standard deviation better
($a_c=-1$): the worse cell has $e^{0.7}=2.0$ times the odds of losing each tape. Two reasons for the
split into $s_\alpha a_c$:
\begin{itemize}
  \item \emph{Common random numbers.} $a_c$ is drawn once and kept, and $s_\alpha$ only rescales it.
        So the same cell is the poorly captured one at every value of $s_\alpha$, and switching the
        scatter off ($s_\alpha=0$) changes nothing else in the dataset.
  \item \emph{Mean zero.} Adding a constant to every $\alpha_c$ and subtracting it from every $\beta_z$
        leaves every $\pi_{cz}$ unchanged, so the overall level can sit in only one of the two. It is
        put in the tapes, as $\mu_\beta$ (\S\ref{sec:pbar}), and $\alpha_c$ averages $0$.
\end{itemize}

\paragraph{Drawing the mask.} $X_{cz}=\mathbb 1[V_{cz}<\pi_{cz}]$, with $V_{cz}\sim\Uniform(0,1)$
drawn once per entry. That is a Bernoulli draw with probability $\pi_{cz}$, written so that raising
$\pi_{cz}$ can only \emph{add} missing entries, never move them: masks at a higher dropout level
contain the masks at a lower one. Given $(\alpha,\beta)$, every $X_{cz}$ is independent of every other
and of the tree.

\paragraph{Why this form.} It is the $M_3$ rung of the nested ladder (\texttt{77}), so on rung-1 data
the ladder fits the true model, and any lineage signal it reports is a false positive. The mask covers
the \emph{whole tape}: internal gaps are $0.02\%$ of Park's entries, and the flat per-cell depth profile
rules out 3$'$ truncation (park-compatibility Fig.~3d/e).

\paragraph{2. The tape panel.} A \emph{panel} is one draw of the $k$ tapes' properties --- for each
tape its class $u_z$, recovery offset $\beta_z$ and speed $r_z$ (\S\ref{sec:latent}). It is drawn once
per simulated experiment and shared by every cell of every tree in it, because all Park clones descend
from one engineered line and carry the same integrations. Tape $z$ is therefore hard to read, or slow,
in every cell and every clone alike; what varies between cells is $\alpha_c$ and the edits.

\paragraph{3. Depth-0 censoring.} A tape that was captured but never edited is indistinguishable, in
Park's pipeline, from a tape that was not captured: \texttt{23\_dropout\_matrix.py} codes any tape whose
first site is empty as \texttt{ABSENT}. The simulator applies the same rule:
\begin{equation}
  Y_{cz}=(1-X_{cz})\,\mathbb 1[D_{cz}\ge1],\qquad O_{cz}=1-Y_{cz}.
  \label{eq:observe1}
\end{equation}
$\mathbb 1[\cdot]$ is $1$ when the condition holds and $0$ otherwise, and $1-X_{cz}$ is $1$ when the tape
was captured. The product of two $0/1$ factors is $1$ only if both are, so it is a logical AND: a tape
is \emph{read} ($Y=1$) only if it was captured \emph{and} edited at least once. The four cases:
\begin{center}\small
\begin{tabular}{@{}lcc@{}}
\toprule
 & edited, $D_{cz}\ge1$ & unedited, $D_{cz}=0$ \\
\midrule
captured, $X_{cz}=0$     & read, $Y=1$     & \textbf{missing}, $Y=0$ (censored) \\
not captured, $X_{cz}=1$ & missing, $Y=0$  & missing, $Y=0$ \\
\bottomrule
\end{tabular}
\end{center}
Only the top-right cell is new: a captured tape lost because it carries nothing to read. Censoring is
on by default and can be switched off; the true $X$ and $D$ are always stored, so every simulated
missing entry knows why it is missing.

\paragraph{4. Cell filters,} applied when an arm is exported, as Park's pipeline does. Let
\begin{equation}
  R_c=\sum_{z=1}^kY_{cz}
  \label{eq:Rc}
\end{equation}
be the number of tapes read in cell $c$. Two filters act on it:
\begin{itemize}
  \item \emph{QC} keeps a cell only if $R_c\ge R_{\min}$ ($20$ or $100$ of $166$ in Park).
  \item \emph{ClonalBC retention} keeps a cell's clone barcode with probability $\psi(R_c)$, a step
        function given as a table: $\psi(R)=p_j$ for $\ell_j\le R<\ell_{j+1}$. In the \textsc{filter}
        fixture of \S\ref{sec:drop1val} it is $0.3$ below $12$ tapes, $0.5$ from $12$, $0.8$ from $20$
        and $0.95$ from $25$; on Park's Mouse1 it runs from $36\%$ at $20$--$40$ tapes to $83\%$ at
        $140$ or more. The draw is $\mathbb 1[W_c<\psi(R_c)]$ with $W_c$ uniform, and a cell that
        fails gets clone label $-1$, which is how \texttt{23}, \texttt{44} and \texttt{77} already
        exclude it.
\end{itemize}
Both filters keep cells \emph{because} they show many tapes, so the cells that survive are selected
for low missingness. Park's data went through this selection, so Park-comparable output has to as
well; it is also why the closed forms of \S\ref{sec:drop1cf} are scored before the filters.

\paragraph{Deliberately absent from this rung:} heritable loss (rung~2), clone-level capture (cells
of one clone sharing capture on every tape --- a cell-axis effect not yet callable as biology), the
\emph{trans} limb of silencing (a silenced integration's symbol disappearing from other tapes), and
the site-6 mechanism.

\subsection{The closed-locus latent: why recovery and speed move together}\label{sec:latent}

\paragraph{The problem it solves.} In Park a minority of tapes are both rarely read \emph{and}
shallow. On Mouse1 the $15\%$ of tapes with recovery below $0.3$ average $3.58$ filled sites against
$4.86$ for the rest, and among the rest recovery and depth are barely related (correlation
$+0.119$). That is a \emph{step}, not a slope. The likely mechanism is a closed integration site,
which is both poorly transcribed (Park reads tapes from RNA) and poorly reached by the prime editor.
Drawing recovery and speed independently would delete that link, and with it the reason editing
and dropout cannot be fitted separately. It would also understate censoring, since slow tapes are
the ones most often left unedited.

\paragraph{The construction.} One hidden switch per tape drives both quantities:
\begin{equation}
  u_z=\mathbb 1[U_z<\phi],\qquad
  \beta_z=\mu_\beta+\Delta_\beta u_z+s_\beta\varepsilon_z,\qquad
  r_z=\bar r_0\,r_{\rm cl}^{\,u_z}\,e^{s_rv_z-s_r^2/2},\qquad
  \bar r_0=\frac{1}{1-\phi+\phi\,r_{\rm cl}} .
  \label{eq:panel}
\end{equation}

\paragraph{How to read it: draws and parameters.} Every symbol in eq.~\eqref{eq:panel} is one of two
kinds. Three are \emph{standard random draws}, made once per tape and kept: $U_z\sim\Uniform(0,1)$ and
$\varepsilon_z,v_z\sim\mathcal N(0,1)$, all independent. They carry no units and no biology; they say
where tape $z$ falls among tapes. The rest are \emph{parameters}, set before a run, which turn those
draws into a class, a recovery and a speed. Changing a parameter rescales the same draws rather than
redrawing them, so a tape keeps its identity across settings.
\begin{center}\small
\begin{tabular}{@{}llp{0.42\textwidth}ll@{}}
\toprule
draw & parameter & what the parameter controls & units & off when \\
\midrule
$U_z$           & $\phi$          & fraction of tapes that are closed               & fraction      & $\phi=0$ \\
---             & $\Delta_\beta$  & how much harder a closed tape is to read        & logits        & $0$ \\
$\varepsilon_z$ & $s_\beta$       & scatter in readability among tapes of one class & logits        & $0$ \\
---             & $r_{\rm cl}$    & how fast a closed tape edits, relative to an open one & ratio   & $1$ \\
$v_z$           & $s_r$           & scatter in speed among tapes of one class       & log-speed     & $0$ \\
---             & $\mu_\beta$     & overall readability; \emph{solved}, not set (\S\ref{sec:pbar}) & logits & --- \\
\bottomrule
\end{tabular}
\end{center}
With every parameter at its ``off'' value, all tapes are identical: equally readable, all at speed
$1$. Each line of eq.~\eqref{eq:panel} now in turn.

\paragraph{The class.} $u_z=\mathbb 1[U_z<\phi]$ is $1$ (closed) with probability $\phi$. Drawing it
through a uniform rather than as a fresh coin flip makes the closed set \emph{nested}: a tape closed at
$\phi=0.15$ is still closed at $\phi=0.3$, so comparisons across $\phi$ are not Monte Carlo noise.

\paragraph{Recovery,} $\beta_z=\mu_\beta+\Delta_\beta u_z+s_\beta\varepsilon_z$, has three terms:
\begin{itemize}
  \item $\mu_\beta$, the level every tape starts from;
  \item $\Delta_\beta u_z$, a penalty added only to closed tapes. At $\Delta_\beta=2$, a tape that would
        be missing $30\%$ of the time if open is missing $76\%$ of the time if closed;
  \item $s_\beta\varepsilon_z$, the \emph{recovery scatter}: tapes of the same class still differ in
        how well they are read, for reasons the class does not capture (sequence context, amplicon
        efficiency). $\varepsilon_z$ says how many standard deviations tape $z$ sits from its class's
        typical value, and $s_\beta$ is that standard deviation in logits. At $s_\beta=1$, among open
        tapes whose typical missingness is $30\%$, a tape one standard deviation worse is missing $54\%$
        of the time and one a standard deviation better $14\%$.
\end{itemize}
Together: given its class, a tape's recovery offset is normal,
$\beta_z\mid u_z\sim\mathcal N(\mu_\beta+\Delta_\beta u_z,\;s_\beta^2)$ --- two bell curves, the closed
one shifted by $\Delta_\beta$. \emph{The scatter decides whether the two classes look like a step or a
slope.} The curves sit $\Delta_\beta/s_\beta$ standard deviations apart. At $s_\beta=0$ they are two
spikes, a pure step. At the validation's $\Delta_\beta=2$, $s_\beta=1$ they are two standard deviations
apart: still distinct (only $2.3\%$ of open tapes are worse read than the typical closed tape) but
overlapping at the edges. As $s_\beta$ grows past $\Delta_\beta$ the two merge into one broad spread and
the step becomes a slope. So the recovery-$0.3$ ``threshold'' seen in Park is \emph{not} a simulator
parameter: it enters only when choosing $(\phi,\Delta_\beta,s_\beta)$ to resemble Park.

\paragraph{Speed,} $r_z=\bar r_0\,r_{\rm cl}^{\,u_z}\,e^{s_rv_z-s_r^2/2}$, is easier to read after
taking logs, where it has the same three-term shape as recovery:
\begin{equation}
  \ln r_z=\underbrace{\ln\bar r_0}_{\text{level}}
  +\underbrace{u_z\ln r_{\rm cl}}_{\text{closed slowdown}}
  +\underbrace{s_rv_z}_{\text{speed scatter}}
  -\underbrace{s_r^2/2}_{\text{mean correction}} .
  \label{eq:lnr}
\end{equation}
Speed is modelled on the log scale because it must stay positive and because tape speeds differ by
\emph{factors}. Term by term:
\begin{itemize}
  \item $r_{\rm cl}^{\,u_z}$ is $1$ for an open tape and $r_{\rm cl}$ for a closed one, so it switches
        the slowdown on only for closed tapes. At $r_{\rm cl}=0.6$ a closed tape edits at $60\%$ of an
        open tape's rate.
  \item $e^{s_rv_z}$ is the \emph{speed scatter}: tapes of one class still edit at different speeds.
        $s_r$ is the standard deviation of $\ln r_z$ within a class. At $s_r=0.3$ a tape one standard
        deviation fast edits $1.29$ times as fast as the class typical, and one a standard deviation slow
        $0.71$ times.
  \item $e^{-s_r^2/2}$ corrects the mean. For a standard normal $v$, $\E[e^{s_rv}]=e^{s_r^2/2}>1$:
        multiplying by $e^{s_rv}$ on its own would make tapes faster \emph{on average}, because fast
        draws are stretched more than slow ones are shrunk. Dividing by $e^{s_r^2/2}$ makes the factor
        average exactly $1$, so the scatter spreads speeds without changing their mean.
  \item $\bar r_0$ is the speed of a typical open tape, chosen so that the average tape has speed $1$.
        Averaging over the class and the scatter,
        \begin{equation}
          \E[r_z]=\bar r_0\big[(1-\phi)\cdot1+\phi\,r_{\rm cl}\big]\cdot\E\big[e^{s_rv-s_r^2/2}\big]
          =\bar r_0\,(1-\phi+\phi r_{\rm cl})\cdot1=1 .
          \label{eq:Er}
        \end{equation}
        At $\phi=0.15$, $r_{\rm cl}=0.6$: $\bar r_0=1/0.94=1.064$, so open tapes run at $1.064$ and closed
        ones at $0.638$. Why it matters: $\LamT$ is defined (eq.~\ref{eq:rate}) as the edits an
        \emph{average} tape receives. Without $\bar r_0$, raising $\phi$ would quietly lower the total
        amount of editing, and an effect of closed tapes would be confounded with an experiment that
        simply edits less.
\end{itemize}

\paragraph{The coupling, and what the scatter does to it.} $\varepsilon_z$ and $v_z$ are independent,
so \emph{within} a class recovery and speed are unrelated. The class is the only link. To see its size,
note that in eqs.~\eqref{eq:panel} and \eqref{eq:lnr} the only random terms are $u_z$ (shared by both),
$\varepsilon_z$ (recovery only) and $v_z$ (speed only). Covariances between independent terms vanish,
leaving the $u_z$ terms:
\begin{equation}
  \mathrm{Cov}(\beta_z,\ln r_z)=\mathrm{Cov}(\Delta_\beta u_z,\;u_z\ln r_{\rm cl})
  =\Delta_\beta\ln r_{\rm cl}\;\mathrm{Var}(u_z)
  =\Delta_\beta\ln r_{\rm cl}\;\phi(1-\phi) .
  \label{eq:latentcov}
\end{equation}
It is negative because $\ln r_{\rm cl}<0$: harder to read goes with slower. The scatter terms do not
enter the covariance, but they do enter the variances, so they dilute the \emph{correlation}:
\begin{equation}
  \mathrm{Corr}(\beta_z,\ln r_z)=\frac{\Delta_\beta\ln r_{\rm cl}\;\phi(1-\phi)}
  {\sqrt{\big[\Delta_\beta^2\phi(1-\phi)+s_\beta^2\big]\big[(\ln r_{\rm cl})^2\phi(1-\phi)+s_r^2\big]}} .
  \label{eq:latentcorr}
\end{equation}
With no scatter ($s_\beta=s_r=0$) this is exactly $-1$: readability and speed are both functions of
the class alone. At the validation's \textsc{cis} values ($\phi=0.15$, $\Delta_\beta=2$,
$r_{\rm cl}=0.6$, $s_\beta=1$, $s_r=0.3$) the covariance is $-0.130$ and the correlation $-0.30$
(checked by Monte Carlo, $4\times10^6$ tapes: $-0.302$). That is the role of the two scatter
parameters: \emph{the class sets how strongly recovery and speed are tied, and the scatter sets how much
independent variation is laid on top}, which is what makes the tie look loose when tapes are viewed
one at a time.

\paragraph{What the scatter is not.} $s_\alpha$, $s_\beta$ and $s_r$ describe fixed differences between
cells or between tapes, drawn once. They are not measurement noise: tape $z$'s $\beta_z$ and $r_z$ are
the same in every cell and every tree of the experiment, and cell $c$'s $\alpha_c$ is the same on every
tape. The randomness from entry to entry is in $X_{cz}$ alone.

\paragraph{A consequence worth knowing.} Even with no closed tapes, the share of unedited tapes exceeds
$e^{-\LamT}$. Speed scatter alone does it: a tape is unedited with probability $e^{-r_z\LamT}$, which is
curved in $r_z$, so by Jensen's inequality $\E[e^{-r\LamT}]>e^{-\E[r]\LamT}=e^{-\LamT}$ --- slow tapes
add more unedited entries than fast ones remove. At $\LamT=2.9$, $s_r=0.3$ the expectation is $0.075$
against $0.055$; the library's five realised panels give $0.079$. Expected, not a finding.

\subsection{\texorpdfstring{Setting the level: the knob is $\pbar$}{Setting the level: the knob is pbar}}\label{sec:pbar}

\paragraph{The idea.} The quantity a user wants to set is how much technical dropout there is, not a
logit offset. So the swept knob is the target \emph{mean technical missingness} $\pbar$, and the level
$\mu_\beta$ is solved from it. $\pbar$ counts only the mask $X$; observed missingness $O$ is higher,
because censoring adds the unedited tapes on top.

\paragraph{The definition.} For tape $z$, the expected missingness over the population of cells is
\begin{equation}
  \pbar_z=\E_\alpha\big[\sigma(\alpha+\beta_z)\big]
  =\int_{-\infty}^{\infty}\sigma(\beta_z+s_\alpha a)\,\frac{e^{-a^2/2}}{\sqrt{2\pi}}\,da ,
  \qquad
  \pbar=\frac1k\sum_{z=1}^k\pbar_z .
  \label{eq:pbar}
\end{equation}
$\E_\alpha$ averages over the \emph{distribution} of cell offsets, $\alpha=s_\alpha a$ with $a$ standard
normal, not over the cells a tree happens to contain; $\pbar$ then averages over the panel's $k$ tapes.

\paragraph{Why the spread has to be averaged over.} One might set $\sigma(\beta_z)$ to the target and be
done. That is wrong whenever the cells differ: $\sigma$ is curved, so averaging it over a spread of
$\alpha$ pulls the result toward $0.5$. A tape at $\beta_z=\operatorname{logit}(0.3)$ is missing in
$30\%$ of \emph{average} cells but in $31.7\%$ of cells \emph{overall} at $s_\alpha=0.7$. To hit a target
below $0.5$, the logits therefore have to sit lower than the naive value.

\paragraph{The solve.} Write each tape's offset from the level as $b_z=\Delta_\beta u_z+s_\beta
\varepsilon_z$, known once the panel is drawn, so that $\beta_z=\mu_\beta+b_z$. Then $\mu_\beta$ is the
root of
\begin{equation}
  F(\mu_\beta)=\frac1k\sum_{z=1}^k\int\sigma(\mu_\beta+b_z+s_\alpha a)\,
  \frac{e^{-a^2/2}}{\sqrt{2\pi}}\,da=\pbar .
  \label{eq:pbarsolve}
\end{equation}
$F$ rises steadily from $0$ to $1$ as $\mu_\beta$ runs from $-\infty$ to $+\infty$, so the root exists
and is unique. The integral is computed by 80-point Gauss--Hermite quadrature (a weighted sum over fixed
nodes, exact for smooth integrands against a normal weight to far beyond the precision needed) and the
root by Brent's method on $[-80,80]$. At $\pbar=0$ the root is $\mu_\beta=-\infty$: no technical
missingness at all.

\paragraph{What follows from solving on the realised panel.}
\begin{itemize}
  \item \emph{$\mu_\beta$ absorbs the panel's own draws.} If a panel happens to draw many hard-to-read
        tapes, $\mu_\beta$ moves down to compensate, so every panel delivers the target on average.
  \item \emph{Raising $\phi$ or $\Delta_\beta$ redistributes missingness rather than adding it.} At fixed
        $\pbar$, more closed tapes means a lower $\mu_\beta$, so the open tapes become \emph{better} read.
        In the validation fixtures, all at $\pbar=0.3$: $\mu_\beta=-1.080$ with no closed tapes
        (\textsc{base}), $-1.176$ with $3$ closed (\textsc{cis}), $-1.667$ with $9$ closed
        (\textsc{slow}). The latent changes \emph{which} tapes are missing, not how much is missing.
  \item \emph{The scatter moves the level too.} With $s_\beta=1$ and $s_\alpha=0.7$, hitting
        $\pbar=0.3$ in expectation over panels needs $\mu_\beta=-1.089$, against the naive
        $\operatorname{logit}(0.3)=-0.847$; with $s_\beta=0$ it needs $-0.937$.
  \item \emph{The target is an expectation.} A finite set of cells realises it with noise: across the
        library the realised technical missingness sits on target on average, with a spread of $0.0046$
        at $n=25$ cells falling to $0.0005$ at $n=1{,}976$.
\end{itemize}

\subsection{Depth-0 censoring makes observed missingness heritable}\label{sec:censor}

This is the one non-obvious property of rung~1, and it is a design result rather than a code check.
Rung~1 contains no silencing, yet \emph{its observed missingness is correlated along the tree},
because ``no edit yet'' is itself inherited.

\paragraph{The derivation.} Take two cells $c,c'$ whose last common ancestor $A$ sits at $t_A$, and
write $W_A=W(t_A)$. On tape $z$:
\begin{enumerate}
  \item A cell is unedited iff its whole root-to-cell path carries no edit. The path is one Poisson
        process in $\Lambda$-time (\S\ref{sec:editcf}, A), so $\Prob(D_c=0)=e^{-\mu_z}$.
  \item Both are unedited iff the \emph{shared} path down to $A$ (expected $\mu_zW_A$ edits) and the
        \emph{two private} paths ($\mu_z(1-W_A)$ each) all carry none. These are three independent
        stretches of the process, so
        \begin{equation}
          \Prob(D_c=0,\,D_{c'}=0)=e^{-\mu_zW_A}\,e^{-2\mu_z(1-W_A)}=e^{-\mu_z(2-W_A)} .
          \label{eq:bothunedited}
        \end{equation}
  \item By inclusion--exclusion, both are edited with probability
        $1-2e^{-\mu_z}+e^{-\mu_z(2-W_A)}$.
  \item The mask is independent of depth, and across cells given $(\alpha,\beta)$, so
        $\Prob(Y_c=1)=(1-\pi_{cz})(1-e^{-\mu_z})$ and
        $\Prob(Y_c=1,Y_{c'}=1)=(1-\pi_{cz})(1-\pi_{c'z})\big(1-2e^{-\mu_z}+e^{-\mu_z(2-W_A)}\big)$.
  \item Averaging over $\alpha$ and subtracting the product of the marginals,
        \begin{equation}
          \boxed{\;
          \mathrm{Cov}(O_{cz},O_{c'z}\mid t_A)=(1-\pbar_z)^2\Big[e^{-\mu_z(2-W_A)}-e^{-2\mu_z}\Big]
          \;\xrightarrow{\ W_A\to1\ }\;(1-\pbar_z)^2\,e^{-\mu_z}\big(1-e^{-\mu_z}\big). \;}
          \label{eq:censorcov}
        \end{equation}
\end{enumerate}
\emph{Reading it.} The no-effect value is $0$, for two cells that split at the founder ($W_A=0$).
The covariance rises the later the pair split, and the sister limit on the right is its ceiling. It
is largest on slow tapes (small $\mu_z$), which by \S\ref{sec:latent} are also the poorly read ones.
In the slowest validation fixture ($\mu=2$, $\pbar=0.3$) the sister covariance is
$0.49\times e^{-2}(1-e^{-2})=0.057$, against $0$ for unrelated cells. At the mouse arms' $\mu\approx5.5$
it is at most $0.004$.

\paragraph{Why it matters for the silencing results.} The heritable-silencing statistics find
relatives on one half of the tapes and measure dropout on the other half, which defeats the
circularity of reading relatedness from the same edits whose readability dropout controls. It does
\emph{not} defeat this, because the shared state lives \emph{within} one tape and the relatives are
found correctly. So depth-0 censoring alone can make dropout look heritable, and it can only push the
silencing readings up. It bears most on the arm with the fewest edits per tape, Pre-TX (variogram
top bin $+0.2618$; ladder $+4.59$ nats per cell, $Q=7.1\%$; unanimity $9.77\times$ against $1.22\times$;
$2{,}019$ events, $2.46\%$ of missing entries), less on Subclone, and essentially not on the mice. Its
size on Park is \textbf{unknown}, and depends on $\Lambda_{\rm pre}$: a tape already edited before the
clone founder cannot be unedited within the clone. Measuring it is the job of the rung-1 negative
control (censoring off against on), which has not yet been run.

\paragraph{Status.} Verified twice: against a two-cell Monte Carlo ($4$ settings, $4\times10^6$ draws
each, $|z|\le0.87$, technical-only covariance $0$ to $10^{-4}$), and in the generator
(\S\ref{sec:drop1val}).

\subsection{What is exported}\label{sec:export}

The point of rung~1 is that simulated data go through \emph{the same code} as Park's. So the
generator writes the two cache files the heritable-silencing scripts read, in their formats:
\begin{itemize}
  \item \texttt{dropout\_matrix\_*.npz}: per cell, the clone label (or $-1$), $R_c$, and per tape a
        term code --- \texttt{ABSENT} if $Y_{cz}=0$, \texttt{COMPLETE} if the tape is full
        ($D_{cz}=N$), \texttt{UNEDITED} otherwise (``read, not yet full'', in Park's naming);
  \item \texttt{prefix\_codes6\_*.npz}: for each depth $j=1,\dots,6$, an integer code for the first
        $j$ symbols of each tape, $-1$ past the tape's depth or where the tape is unobserved; two cells
        share a code at depth $j$ iff their first $j$ symbols agree. Its rows are the cells that pass
        QC and carry a clone label, in dropout-matrix order (\texttt{44}'s alignment rule).
\end{itemize}
\texttt{73}, \texttt{74} and \texttt{77} gained two arguments, \texttt{--resdir} and \texttt{--thr}
(input and output paths only, no logic), and a simulated arm loads and runs through all three.

\subsection{The closed forms rung 1 is checked against}\label{sec:drop1cf}

\begin{description}
  \item[F0 --- the panel law] (unit: one panel, $20{,}000$ panels):
        $\E[u]=\phi$; $\E[r]=1$; the closed share of speed $\E[ru]=\phi\bar r_0r_{\rm cl}$;
        $\E[\beta-\mu_\beta]=\Delta_\beta\phi$; and
        $\E[(\beta-\mu_\beta)\ln r]=\Delta_\beta\phi(\ln\bar r_0+\ln r_{\rm cl}-s_r^2/2)$.
  \item[F1 --- the mask:] $\Prob(X=1\mid\eta)=\sigma(\eta)$, $\eta=\alpha_c+\beta_z$.
  \item[F2 --- observed missingness per tape:] $\Prob(O_{cz}=1)=1-(1-\pi_{cz})(1-e^{-\mu_z})$.
  \item[F3 --- depth given read:]
        $\Prob(D=d\mid Y=1,z)=\Prob\big(\min(\Poisson(\mu_z),N)=d\big)/(1-e^{-\mu_z})$, $d=1,\dots,N$,
        by locus class. The mask is blind to depth, so closed tapes look shallower through $r_z$ alone.
  \item[F4t --- the mask never reads the tree:] $\E[X_{cz}X_{c'z}]=\pi_{cz}\pi_{c'z}$ at every split
        time, conditional on the realised $\alpha$.
  \item[F4o --- inherited censoring:] the pair form of eq.~\eqref{eq:censorcov}, conditional on the
        realised $\alpha$ (step 5 before averaging).
  \item[F5 --- retention:] $\Prob(\text{kept}\mid R_c)=\mathbb 1[R_c\ge R_{\min}]\,\psi(R_c)$.
\end{description}
F1--F4 are scored \emph{before} the filters. After them the marginals are no longer
$\sigma(\eta)$, because keeping a cell conditions on $R_c$ and $R_c$ depends on $X$. That selection is
what Park-comparable output has to have, not a defect. F4 pairs are taken one per (internal node,
tape), one cell from each side of the split, never random cell pairs --- the lesson of
\S\ref{sec:editval}.

% ===========================================================================
\section{Validation of rung 1}\label{sec:drop1val}
% ===========================================================================

Validation (\texttt{18\_validate\_dropout.py}) asks whether the code generates the model of
\S\ref{sec:drop1}. A failure is a bug, not a finding. Every entry is a deviation from a closed form of
\S\ref{sec:drop1cf} in standard errors, eq.~\eqref{eq:z}, with the standard error taken across trees.

\paragraph{Four fixtures.} $n=210$, $k=30$, $N=6$, the same 200 library trees as
\S\ref{sec:editval} ($\capt$ = $0.0005$, $0.002$, $0.02$, $0.1$, $0.25$; $\turn$ = $0.3$, $0.7$), and one mask
setting throughout: $\pbar=0.3$, $s_\beta=1$, $s_\alpha=0.7$. The fixtures are \emph{test cases, not
science}: each is chosen so that one check has something to detect.
\begin{center}\small
\begin{tabular}{@{}lp{0.70\textwidth}@{}}
\toprule
name & what it sets \\
\midrule
\textsc{base}   & constant rate, $\LamT=5.5$; no closed tapes ($\phi=0$), no speed scatter.
                  Censoring is rare: $0.4\%$ of tapes unedited \\
\textsc{cis}    & \textsc{base} $+$ the latent: $\phi=0.15$, $r_{\rm cl}=0.6$, $\Delta_\beta=2$, $s_r=0.3$
                  (3 closed tapes in the realised panel) \\
\textsc{slow}   & the \textsc{cis} panel law at $\LamT=2$ with the \textsc{tv} pause-and-burst rate of
                  \S\ref{sec:editval} (9 closed tapes); $2.5$--$45\%$ of a tape's entries unedited,
                  so censoring is large \\
\textsc{filter} & \textsc{base} $+$ $R_{\min}=12$ of $30$ and a four-step $\psi$:
                  $0.3/0.5/0.8/0.95$ from $R_c=0/12/20/25$ \\
\bottomrule
\end{tabular}
\end{center}

\paragraph{The reading rule, fixed before the run.} No-effect value $|z|=0$. \textbf{Pass} means the
worst $|z|$ over every verdict cell is at most $3.0$ \emph{and} there are no violations of the hard
assertions. With about $220$ cells, one $|z|$ above $3$ is expected about $0.6$ times by chance, so a
single excursion is re-run at a fresh seed before it counts; two, or a sign pattern within one check,
is a defect. \emph{Decisive failure:} F4t not flat (the mask reads the tree), any encoding
assertion, or F4o failing while F2 passes. That last would mean depth~0 is not inherited the way
\S\ref{sec:editcf} says, or that \texttt{23}'s \texttt{ABSENT} rule is not what it was read to be ---
re-derive rather than patch.

\begin{center}\small
\begin{tabular}{@{}l p{0.56\textwidth} r r@{}}
\toprule
check & what it tests & cells & worst $|z|$ \\
\midrule
F0  & the panel law, two laws (\textsc{cis}, and a stronger $\phi=0.3$, $r_{\rm cl}=0.4$)      & 10  & 1.92 \\
F1  & technical missingness against $\sigma(\eta)$, 10 bins $\times$ 4 fixtures                & 40  & 1.84 \\
F2  & observed missingness per tape                                                             & 120 & 2.98 \\
F3  & depth given read, by locus class                                                          & 36  & 2.06 \\
F4t & both relatives technically missing, one pooled cell per fixture                           & 4   & 0.91 \\
F4o & both relatives observed missing, pooled, \textsc{cis} and \textsc{slow} only              & 2   & 1.17 \\
F5  & retention by $R_c$                                                                        & 4   & 1.87 \\
H   & $Y=(1-X)\mathbb 1[D\ge1]$ exactly; $Y=0$ wherever $D=0$; term codes, depths, $R_c$ and clone
      labels in the export; prefix codes determined by the symbols and in bijection with them;
      the solve for $\mu_\beta$; replay bit-identical
      & --- & $0$ of $2.7{\times}10^{7}$ \\
\bottomrule
\end{tabular}
\end{center}

\paragraph{Result: pass, after a disclosed failure.} The verdict above is the third run. The history
matters, because the failure was in the check set and not the generator.
\begin{itemize}
  \item \emph{Seed 29092026:} one trip (F2 \textsc{base}, tape 11, $z=3.16$; $0.78$ expected by
        chance), so re-run by rule.
  \item \emph{Seed 30092026:} two trips, F4t \textsc{base} bin 3 ($-3.54$) and F4o \textsc{base}
        bin 3 ($-3.38$): \textsc{fail} by rule. Zero hard violations at both seeds.
  \item \emph{Diagnosis.} The two trips were one fluctuation counted twice. Where censoring is rare,
        observed missingness \emph{is} technical missingness, so F4o and F4t measure the same thing:
        their $z$ values correlated $+0.98$ to $+1.00$ across bins in \textsc{base} and
        \textsc{filter}. And F4t cannot be biased by construction --- independent uniforms, cells
        from disjoint subtrees.
  \item \emph{Follow-up} (\texttt{21\_f4\_followup.py}, approved separately): each F4 check pooled
        into one pre-declared cell per fixture, on $840$ trees (all $42$ stored $(\capt,\turn)$ cells
        at $n=210$) $\times$ $5$ fresh panels, plus a new check F4d, ``both relatives unedited'',
        against eq.~\eqref{eq:bothunedited} --- censoring alone, no mask. Rule: $|z|\le3$ on all six
        cells means no defect.
\end{itemize}

\begin{center}\small
\emph{Reading rule:} each row is a probability for one (internal node, tape) comparison, pooled over
$840$ trees and $5$ panels; obs $-$ pred near $0$ and $|z|\le3$ means the generator matches the closed
form.\\[3pt]
\begin{tabular}{@{}llrrrr@{}}
\toprule
fixture & check & observed & predicted & obs $-$ pred & $z$ \\
\midrule
\textsc{slow} & F4d both unedited            & 0.09848 & 0.09863 & $-0.00016$ & $-0.64$ \\
\textsc{slow} & F4t both technically missing & 0.12425 & 0.12422 & $+0.00003$ & $+0.42$ \\
\textsc{slow} & F4o both observed missing    & 0.22661 & 0.22670 & $-0.00009$ & $-0.36$ \\
\textsc{cis}  & F4d                          & 0.00505 & 0.00507 & $-0.00002$ & $-0.49$ \\
\textsc{cis}  & F4t                          & 0.12851 & 0.12849 & $+0.00002$ & $+0.31$ \\
\textsc{cis}  & F4o                          & 0.13545 & 0.13549 & $-0.00003$ & $-0.37$ \\
\bottomrule
\end{tabular}
\end{center}
\noindent No defect: worst pooled $|z|=0.64$, with power to spare (the \textsc{slow} F4o standard error
is $0.00026$, so the earlier seed's excess would have been about $10$ standard errors). The scoring
was then amended, in the validation code only with the generator untouched: \textbf{each F4 check
counts through one pooled cell per fixture}, its ten $W$-bins kept as secondary output, and
\textbf{F4o counts only where censoring is not negligible} (\textsc{cis}, \textsc{slow}). Re-run at
seed 2102026: $216$ verdict cells, none excluded, worst $|z|=2.98$ (F2 \textsc{cis}), none above $3$,
zero hard violations, split-null spread $0.95$. \emph{Disclosed:} one secondary bin (F4t
\textsc{cis}) reached $3.10$, which would have tripped under the old scoring; its pooled cell is
$-0.54$.

\paragraph{The check with teeth.} Scoring relatives' observed missingness against the
\emph{independence} prediction --- what it would be if censoring were not inherited --- misses by
$z=26.8$ pooled on \textsc{slow}. Eq.~\eqref{eq:censorcov} is real in the generator, and F4o sees it.

\paragraph{Two lessons, both general.} \emph{Bins of one check share their trees}, so a sign pattern
across bins tracks the pooled value and is not independent evidence. And \emph{two statistics that
coincide in a regime must not both count toward a verdict there.} This extends the lesson of
\S\ref{sec:editval}.

\paragraph{Cost.} $0.06$~s per tree at $n=210$, $k=30$ for decoration, mask and censoring together ---
the mask adds little to the editing layer's $0.05$~s; the whole validation runs in $76$~s on one core.

\paragraph{What passing does not establish.} F1--F3 are marginals: a mask aligned with lineage would
pass them, and only F4t and F4o see the joint. Rung~1 is a model, and a Park-shaped run of it inherits
a \textbf{double count}: Park's measured tape offsets $\hat\beta_z$ already contain depth-0 censoring
and heritable loss, so feeding them in as the \emph{technical} rate counts both twice. No Park
parameter enters this build for that reason. And editing and dropout still have to be fitted jointly
before any Park-calibrated claim.

\subsection{The rung-1 library}\label{sec:drop1lib}

Generated only after the pass, and built in three layers so that expensive work is reused:
\begin{enumerate}
  \item \textbf{Trees:} the existing tree library, no new trees --- $n$ = $25$, $64$, $210$, $626$, $1{,}976$;
        $\capt$ = $0.002$, $0.1$, $0.25$; $\turn$ = $0.3$, $0.7$; $20$ trees each.
  \item \textbf{Panel and editing:} for each $\LamT\in\{2.9,5.5\}$ and panel type \emph{none}
        ($\phi=0$), \emph{park} ($\phi=0.15$, $r_{\rm cl}=0.6$) or \emph{strong} ($\phi=0.3$,
        $r_{\rm cl}=0.4$), all with $s_r=0.3$ and $\Delta_\beta=2$: five panel replicates, each shared
        by every tree, and one decoration of each tree. $k=30$, $N=6$, constant rate, lineage channel
        only (Park's pooled mouse $\xi^L$), tapes empty at the founder.
  \item \textbf{Masks:} $30$ per decorated tree, $\pbar\in\{0,0.1,0.25,0.44,0.6\}$ $\times$
        $s_\beta\in\{0,1,2\}$ $\times$ $s_\alpha\in\{0,0.7\}$, from per-tree common random numbers, so
        masks are nested as $\pbar$ rises.
\end{enumerate}
The truth ($D$, symbols, $X$, the panel) is stored; censoring and filters are applied when a dataset
is read, so ``censoring off'' costs nothing. \emph{Result:} $180$ files, $27{,}000$ datasets, $1.9$~GB
(regenerable from stored seeds), zero assertion violations. Nothing is concluded from it; its manifest
checks only that realised technical missingness sits on its target (mean difference $\approx0$; spread
$0.0046$ at $n=25$ falling to $0.0005$ at $n=1{,}976$), that $\pbar=0$ datasets carry exactly no
technical dropout, and that the closed set is nested (\emph{park} $2$--$4$, \emph{strong} $7$--$12$
closed tapes of $30$).

% ===========================================================================
\section{The dropout layer, rung 2: heritable silencing}\label{sec:drop2}
% ===========================================================================

\paragraph{The idea.} An integration site can be silenced in a cell, and the cell's descendants
inherit the silence. The tape then reads as missing in that whole subclade, for good. This is row~A9
of the assumption table, and the earlier work on Park showed it happens: related cells lose the same
tapes, and the loss is discrete rather than graded. Rung~2 adds it on top of rung~1, in its simplest
form: \emph{complete} Dollo loss --- a tape silenced in a cell is unreadable in every descendant, never
returns, and no descendant escapes.

\paragraph{Four design decisions.} (1) The knob is the \emph{target silenced fraction} $\Hbar$, and
the rate is solved from it. (2) Silencing happens per unit \emph{time}, with its timing set by a shape
$\zeta_0(t)$ swept as its own axis. (3) There is no silencing before the clone founder: one tree is
one clone, and its tapes are intact at the founder. (4) Every tape has the same hazard, since
tape-to-tape differences in propensity cannot be read off Park as it stands (\S\ref{sec:drop2val}).

\subsection*{Notation}

As \S\ref{sec:drop1}, plus the following. Three letters that would be natural here are avoided
because they are taken: $\gamma$ (the clone $\times$ tape term fitted by \texttt{77}), $G$ and
$\omega_g$ (the editing epochs, \S\ref{sec:piecewise}) and $h$ (eq.~\eqref{eq:G}).

\begin{center}\small
\begin{tabular}{@{}lp{0.55\textwidth}l@{}}
\toprule
symbol & meaning & units / range \\
\midrule
$e$, ${\rm pa}(e)$ & a branch --- the edge \emph{into} node $e$, on $[t_e^0,t_e^1]$; the stem is
      $[0,t_{\rm root}]$ --- and its parent branch & --- \\
$n_e$ & number of cells below branch $e$: its clade size & cells, $1$--$n$ \\
$\zeta(t)$ & silencing hazard, per tape, per lineage & events / tape /\\ & & lineage / experiment \\
$\GamT$ & expected silencing events per tape along one founder-to-cell lineage over the whole
      experiment, counted as if every event could strike & events / tape \\
$\zeta_0(t)$ & shape of the hazard in time, $\int_0^1\zeta_0=1$; piecewise constant on its own knots &
      dimensionless \\
$W_\zeta(t)$ & $\int_0^t\zeta_0$: share of the experiment's silencing hazard spent by $t$ & $[0,1]$ \\
$\Gamma(t)$, $\Gamma_e$ & $\GamT W_\zeta(t)$, a lineage's exposure by $t$; and the hazard carried by
      branch $e$ & events / tape \\
$\Hbar$ & target silenced fraction: expected share of cell $\times$ tape entries silenced at harvest,
      before the mask, censoring and filters & probability \\
$E_{e,z}$, $L_{e,z}$, $t^*_{e,z}$ & the $\Exp(1)$ draw for branch $e$, tape $z$; $1$ if a silencing
      event fires on that branch; and the time of the first one. \emph{$L_{e,z}$ is not the lineage
      count $L(t)$ of Fig.~S3} & ---, $\{0,1\}$, time \\
$H_{cz}$ & $1$ if tape $z$ is silenced in cell $c$ (also $H_{e,z}$ along branches) & $\{0,1\}$ \\
$V_{e,z}$ & $1$ if branch $e$ carries the \emph{first} loss on its lineage: a visible event &
      $\{0,1\}$ \\
\bottomrule
\end{tabular}
\end{center}

\subsection{The loss process}\label{sec:loss}

Silencing strikes each tape along each lineage as a Poisson process with hazard
\begin{equation}
  \zeta(t)=\GamT\,\zeta_0(t),\qquad \int_0^1\!\zeta_0(u)\,du=1 .
  \label{eq:zeta}
\end{equation}
$\GamT$ sets \emph{how much} silencing there is and $\zeta_0$ sets \emph{when}, exactly as $\LamT$ and
$\lambda_0$ do for editing (eq.~\ref{eq:rate}); $\zeta_0$ has its own knots, independent of
$\lambda_0$'s. The running total is the silencing analogue of $\Lambda$-time:
\begin{equation}
  \Gamma(t)=\GamT W_\zeta(t),\qquad W_\zeta(t)=\int_0^t\!\zeta_0(u)\,du,\qquad
  \Gamma_e=\Gamma(t_e^1)-\Gamma(t_e^0).
  \label{eq:Gamma}
\end{equation}
$W_\zeta$ runs from $0$ to $1$; $\Gamma_e$ is the hazard branch $e$ carries. As with editing, the tree
only cuts one clock into branches.

\subsection{\texorpdfstring{The knob: the silenced fraction $\Hbar$}{The knob: the silenced fraction}}\label{sec:Hbar}

\paragraph{The solve.} Every cell's founder-to-cell path carries exposure $\Gamma(1)=\GamT$, however
the tree branches. A cell keeps a tape iff no event struck its path, so
\begin{equation}
  \boxed{\;\Hbar=\Prob(H_{cz}=1)=1-e^{-\GamT}\qquad\Longleftrightarrow\qquad\GamT=-\ln(1-\Hbar).\;}
  \label{eq:Hsolve}
\end{equation}
Closed form, no root-finding, and \textbf{independent of the tree}: every tree in a dataset gets the
same $\GamT$. Also independent of $\zeta_0$: moving silencing earlier or later changes \emph{which}
cells share a loss, not how many entries are lost.

\begin{center}\small
\begin{tabular}{@{}lrrrrrrr@{}}
\toprule
target $\Hbar$ (fraction of entries silenced) & 0.01 & 0.02 & 0.05 & 0.10 & 0.20 & 0.25 & 0.33 \\
solved $\GamT$ (events per tape per lineage)  & 0.0101 & 0.0202 & 0.0513 & 0.1054 & 0.2231 & 0.2877 & 0.4005 \\
\bottomrule
\end{tabular}
\end{center}

\paragraph{Expected, not realised.} Silencing is shared by relatives, so the fraction realised in one
tree fluctuates around $\Hbar$. For tape $z$, with $\hat H_z=\frac1n\sum_cH_{cz}$ and $t_A$ the split
time of each pair ($t_A=1$ when $c=c'$),
\begin{equation}
  \mathrm{Var}(\hat H_z)=\frac1{n^2}\sum_{c,c'}e^{-2\GamT}\big(e^{\GamT W_\zeta(t_A)}-1\big).
  \label{eq:Hvar}
\end{equation}
Early silencing hits big clades, so one event moves many entries at once and the realised fraction is
noisier; at the extreme a loss on the stem silences the tape in the whole clone. The solve is
therefore on the expectation, and the realised value is reported. Forcing it to hit $\Hbar$ in every
tree would condition the Poisson process on its outcome and change the law of event sizes, which is
what the downstream statistics read.

\paragraph{$\Hbar$ is not added on top of observed missingness.} The three causes overlap:
\begin{equation}
  \Prob(O_{cz}=1)=1-(1-\pi_{cz})(1-\Hbar)(1-e^{-\mu_z}).
  \label{eq:O3}
\end{equation}

\subsection{Sampling a loss on one branch}\label{sec:losssampler}

The editing sampler of \S\ref{sec:sampler} with one slot instead of $N$. Per branch and tape, draw
$E_{e,z}\sim\Exp(1)$, then
\begin{equation}
  L_{e,z}=\mathbb 1\big[E_{e,z}\le\Gamma_e\big],\qquad
  t^*_{e,z}=W_\zeta^{-1}\!\Big(W_\zeta(t_e^0)+\frac{E_{e,z}}{\GamT}\Big)\quad\text{if }L_{e,z}=1 .
  \label{eq:losssampler}
\end{equation}
\emph{Why it is exact.} The first event on a branch arrives an $\Exp(1)$ amount of hazard after the
branch starts, because the process is memoryless; disjoint branches are independent; and once a tape is
lost, later events change nothing, so only the first is needed. $W_\zeta^{-1}$ is the two-step inverse
of eq.~\eqref{eq:Winv}, so a zero-hazard window needs no special case, just as a zero-rate pause did.

\paragraph{Nesting.} For fixed draws, $L_{e,z}$ only switches on as $\GamT$ rises, and $\GamT$ rises with
$\Hbar$. So the silenced set at a lower $\Hbar$ is a \emph{subset} of the set at a higher one, and a
sweep over $\Hbar$ walks along one path rather than comparing independent draws. Sets are nested
across $\Hbar$ at a fixed $\zeta_0$, \emph{not} across shapes: changing $\zeta_0$ moves hazard between
branches. The loss layer has its own random-number stream, so switching it on never changes the rung-1
draws.

\subsection{Inheritance: the Dollo rule}\label{sec:dollo}

\begin{equation}
  H_{e,z}=H_{{\rm pa}(e),z}\;\vee\;L_{e,z},\qquad H=0\ \text{above the stem.}
  \label{eq:dollo}
\end{equation}
A lineage has lost the tape if it was lost on its own branch or anywhere above it. $H$ can only switch
\emph{on} going down the tree; that one-way property \emph{is} the Dollo character, and it is a hard
assertion. A \emph{visible} event is the first loss on a lineage,
$V_{e,z}=L_{e,z}\wedge\neg H_{{\rm pa}(e),z}$; a loss below an earlier one is invisible. Tapes are intact
at the founder, but a \emph{clone-wide} loss can still happen: it is a loss on the stem, with probability
$1-e^{-\GamT W_\zeta(t_{\rm root})}$. Trees are independent, so no silencing is shared between clones,
even when they share a rung-1 panel.

\subsection{Observation: three independent causes}\label{sec:observe2}

\begin{equation}
  Y_{cz}=(1-X_{cz})\,(1-H_{cz})\,\mathbb 1[D_{cz}\ge1],\qquad O_{cz}=1-Y_{cz}.
  \label{eq:observe2}
\end{equation}
A tape is read only if it was technically captured, \emph{and} not silenced, \emph{and} edited at least
once. Each cause has its own random-number stream and the truth triple $(X,H,D)$ is stored, so every
missing entry knows why it is missing. \textbf{Editing is left untouched}: $D$ comes from the editing
layer whether or not the tape was later silenced. Under complete loss with no reactivation and no
\emph{trans} limb, edits written after silencing are never observed, so the choice cannot be seen in the
data --- and it keeps the rung-1 library's trees, panels and decorations reusable. The export of
\S\ref{sec:export} takes this $Y$ unchanged and writes $Y=0$ as \texttt{ABSENT}.

\subsection{The closed forms rung 2 is checked against}\label{sec:drop2cf}

Each is written in the knob as well, using $e^{-\GamT}=1-\Hbar$.

\paragraph{G1 --- the marginal (tree-blind).} $\Prob(H_{cz}=1)=\Hbar$ for every tape and every cell.
This is eq.~\eqref{eq:Hsolve} and tests the solve. Like the editing layer's check A, it cannot see
\emph{which} lineage a loss landed on.

\paragraph{G2 --- pairs.} Two cells both keep the tape iff three independent pieces carry no event: the
shared path (exposure $\GamT W_\zeta(t_A)$) and the two private paths ($\GamT[1-W_\zeta(t_A)]$ each).
So
\begin{equation}
\begin{gathered}
  \Prob(H_c{=}0,H_{c'}{=}0\mid t_A)=(1-\Hbar)^{\,2-W_\zeta(t_A)},\\[2pt]
  \mathrm{Corr}(H_c,H_{c'}\mid t_A)=\frac{(1-\Hbar)\big[(1-\Hbar)^{-W_\zeta(t_A)}-1\big]}{\Hbar}
  \;\xrightarrow{\ \Hbar\to0\ }\;W_\zeta(t_A).
\end{gathered}
  \label{eq:G2}
\end{equation}
\emph{In words:} when silencing is rare, \textbf{the correlation between two cells' silencing equals the
share of the silencing hazard they lived through together.} The no-effect value is $0$, for a pair split
at the founder. Under constant $\zeta_0$ it is close to linear in split time; under a late-heavy
$\zeta_0$ most pairs share almost nothing and only very recent relatives co-lose. The validation scores
the equivalent ``both silenced'' form, $1-2(1-\Hbar)+(1-\Hbar)^{2-W_\zeta(t_A)}$.

\paragraph{G3 --- where losses land.} This is the check that sees lineage assignment. Branch $e$ carries
a visible event iff nothing fired above it and something fired on it:
\begin{equation}
  \Prob(V_{e,z}=1)=e^{-\Gamma(t_e^0)}\big(1-e^{-\Gamma_e}\big)
  =(1-\Hbar)^{W_\zeta(t_e^0)}\Big[1-(1-\Hbar)^{W_\zeta(t_e^1)-W_\zeta(t_e^0)}\Big],
  \label{eq:G3}
\end{equation}
and the expected number of visible events whose clade size falls in a bin $\mathcal B$, summed over
tapes, is
\begin{equation}
  \E[\mathcal V_{\mathcal B}]=k\sum_{e:\,n_e\in\mathcal B}\Prob(V_{e,z}=1)\qquad\text{events per tree.}
  \label{eq:G3bin}
\end{equation}
At small $\Hbar$, events land on branches in proportion to
the hazard each carries. \emph{Consistency:} $\sum_en_e\Prob(V_{e,z}=1)=n\Hbar$, which is G1, since each
silenced cell's first loss sits on exactly one branch of its path.

\paragraph{The heritable share} --- the quantity to quote alongside $\Hbar$:
\begin{equation}
  \frac{\Hbar_{\ge2}}{\Hbar}=\frac{\sum_{e:\,n_e\ge2}n_e\Prob(V_{e,z}=1)}{n\Hbar}
  \;\xrightarrow{\ \Hbar\to0\ }\;1-\frac1n\sum_c\Big[1-W_\zeta\big(t^0_{e(c)}\big)\Big],
  \label{eq:hershare}
\end{equation}
where $e(c)$ is cell $c$'s terminal branch. It is the fraction of silenced entries that sit in a clade
of at least two cells, which is the only silencing a relatedness-based statistic can ever see.
\emph{Why the limit:} at small $\Hbar$ a cell is silenced on its own terminal branch with probability
about $\GamT[1-W_\zeta(t^0_{e(c)})]$, out of $\GamT$ in total, so the singleton share is the mean of
$1-W_\zeta(t^0_{e(c)})$ over cells --- the average share of a cell's silencing hazard spent on its own
terminal branch. Its floor, $0$, is all silencing in singletons, where it looks exactly like technical
dropout. \textbf{A rung-2 setting is therefore quoted with $\Hbar$, $\zeta_0$ and the heritable share,
never with $\Hbar$ alone.}

\paragraph{G4 --- loss time inside a branch.} Given $L_{e,z}=1$, the probability-integral transform of
the first-event time, truncated to the branch, is uniform:
\begin{equation}
  \frac{1-e^{-\GamT[W_\zeta(t^*)-W_\zeta(t_e^0)]}}{1-e^{-\Gamma_e}}\sim\Uniform(0,1).
  \label{eq:G4}
\end{equation}
It tests $W_\zeta^{-1}$ under a shaped $\zeta_0$. The stored $t^*$ is the ground truth for any later
``when was it silenced'' estimate.

\paragraph{G5 --- the composed view.} $X$, $H$ and $D$ are independent given the panel and the tree, so
\begin{align}
  \Prob(Y_{cz}=1)&=(1-\pi_{cz})(1-\Hbar)(1-e^{-\mu_z}),\nonumber\\
  \Prob(Y_c{=}1,Y_{c'}{=}1\mid t_A)&=(1-\pi_{cz})(1-\pi_{c'z})\cdot(1-\Hbar)^{2-W_\zeta(t_A)}\cdot
  \big(1-2e^{-\mu_z}+e^{-\mu_z(2-W(t_A))}\big).
  \label{eq:G5}
\end{align}
The middle factor is G2 and the last is \S\ref{sec:censor}'s both-edited probability. The product says
the three causes do not interact, which is what makes the stored truth an exact decomposition of the
missingness.

\subsection{Silencing and censoring are one object, mirrored}\label{sec:mirror}

\emph{Derived; each half verified separately; the comparison itself not yet run.} Both are first-event
processes along a lineage. Censoring is missing in the \emph{not-yet-hit} state (no first edit yet);
silencing is missing in the \emph{hit} state (first loss has happened). Since
$\mathrm{Cov}(H,H')=\mathrm{Cov}(1-H,1-H')$, their pair covariances are the same function,
$e^{-\varrho(2-w)}-e^{-2\varrho}$, with $(\varrho,w)=(\mu_z,W)$ for censoring
(eq.~\ref{eq:censorcov}) and $(\GamT,W_\zeta)$ for silencing.

At a matched \emph{rate} they coincide. At a matched missing fraction $m_0$ --- the comparison data
offer --- they come apart. With both hazards constant, so that $w$ is the pair's split time,
\begin{equation}
  \mathrm{Corr}_{\rm cens}(w)=\frac{m_0\big(m_0^{-w}-1\big)}{1-m_0},\qquad
  \mathrm{Corr}_{\rm sil}(w)=\frac{(1-m_0)\big((1-m_0)^{-w}-1\big)}{m_0}.
  \label{eq:mirror}
\end{equation}
\emph{Illustration, not a result:} at $m_0=0.1$ and $w=0.5$ the correlation is $0.24$ for censoring and
$0.49$ for silencing. Rare censoring is \emph{convex} in split time --- only very recent relatives share
it --- while rare silencing rises roughly linearly. This suggests two discriminators for the negative
control: the \emph{convexity} of the dropout variogram, and how deep a cell's observed relatives sit on
the same tape (shallow under censoring, ordinary under silencing). \emph{Caveat:} a late-heavy
$\zeta_0$ also makes silencing convex, so convexity separates the two only together with a stated
$\zeta_0$, which the rung-2 library supplies. No recorded number moves.

% ===========================================================================
\section{Validation of rung 2}\label{sec:drop2val}
% ===========================================================================

\texttt{23\_validate\_silencing.py} reuses rung~1's accumulators and pair sampler, so the scoring is
rung~1's amended scoring: standard errors across trees, pairs one per (internal node, tape), and G2
counted through one pooled cell per fixture with its ten $W_\zeta$-bins secondary.

\paragraph{Five fixtures,} on the same 200 trees ($n=210$, $k=30$), test cases rather than science.
\begin{center}\small
\begin{tabular}{@{}lp{0.72\textwidth}@{}}
\toprule
name & what it sets \\
\midrule
\textsc{off}      & $\Hbar=0$ on the rung-1 \textsc{base} fixture: separability only \\
\textsc{const}    & $\Hbar=0.25$, $\zeta_0$ constant \\
\textsc{early}    & $\Hbar=0.25$, $\zeta_0=3$ on $[0,0.25)$, $0$ on $[0.25,0.6)$, $0.625$ on $[0.6,1]$ ---
                    a zero-hazard window, and big clades \\
\textsc{late}     & $\Hbar=0.25$, $\zeta_0=0$ on $[0,0.7)$, $3.33$ on $[0.7,1]$ --- mostly singletons \\
\textsc{composed} & the rung-1 \textsc{slow} fixture with \textsc{const} silencing: all three causes
                    large at once \\
\bottomrule
\end{tabular}
\end{center}
(The rung-2 \textsc{const} is not the editing layer's \textsc{const} of \S\ref{sec:editval}.)

\paragraph{The reading rule, fixed before the run.} As rung~1: \textbf{pass} = worst $|z|\le3.0$ over
every verdict cell and no hard violations; about $175$ cells, so one excursion ($0.46$ expected) is
re-run at a fresh seed, while two, or a sign pattern within one check, is a defect. \emph{Decisive
failure:} any violation of Dollo monotonicity; \textsc{off} not bit-identical to rung~1; or G3 failing
while G1 passes --- the right amount of silencing in the wrong place. And if G1 and G4 pass but G2 or
G3 fail in a consistent direction, the per-branch hazard and the path total are right but the
inheritance and visibility rules of \S\ref{sec:dollo} are wrong: re-derive before writing more code.

\begin{center}\small
\begin{tabular}{@{}l p{0.58\textwidth} r r@{}}
\toprule
check & what it tests & cells & worst $|z|$ \\
\midrule
G1 & silenced fraction per tape $=\Hbar$ (tests the solve; tree-blind), 3 fixtures         & 90 & 2.79 \\
G2 & relatives both silenced, eq.~\eqref{eq:G2}, one pooled cell per fixture                & 3  & 1.37 \\
G3 & visible events by clade size, eq.~\eqref{eq:G3} --- \emph{where} losses land; size bins
     $1\cdot2\cdot3$--$4\cdot5$--$16\cdot17$--$64\cdot65$--$209\cdot210$ (the stem)           & 18 & 2.22 \\
G4 & loss-time transform inside its branch, deciles, eq.~\eqref{eq:G4}                     & 30 & 1.78 \\
G5 & composed $\Prob(Y=1)$ per tape, and relatives both read, eq.~\eqref{eq:G5}             & 31 & 2.34 \\
H  & \textsc{off} bit-identical to rung~1 in $Y$ and in both exported caches; Dollo monotone; cell $H$
     recomputed independently as ``any loss on the path''; $\sum_en_eV_{e,z}=\sum_cH_{cz}$ exactly;
     $t^*$ inside its branch and never in a zero-hazard window; silenced sets nested across
     $\Hbar=0.02<0.10<0.25<0.33$; $Y$ identity and \texttt{ABSENT} wherever silenced in the export;
     the solve exact; replay identical
   & --- & $0$ of $3.1{\times}10^{7}$ \\
\bottomrule
\end{tabular}
\end{center}

\paragraph{Result: pass, at the second seed.} A 20-tree smoke run (not a verdict) tripped $6$ cells,
worst G1 \textsc{early} at $6.61$, with every hard assertion holding and the pooled means on target.
This was read as a $t$-statistic not yet calibrated at 20 trees --- the situation met in
\S\ref{sec:editval} --- and the full run went ahead with \emph{no change to the checks or scoring}. Seed
1102026: one trip (G1 \textsc{const}, tape 9, $|z|=3.66$; $0.46$ expected), no sign pattern ($17$ of
$30$ tapes negative), so re-run by rule. \textbf{Seed 3102026: pass} --- $172$ verdict cells ($3$
excluded, all G3 \textsc{late} bins of $17$ cells or more, where too few events are expected), worst
$|z|=2.79$, none above $3$, zero hard violations, split-null spread $1.04$.

\paragraph{The checks have teeth.} Two diagnostics outside the verdict, both of which must miss:
\begin{itemize}
  \item relatives' co-silencing scored against \emph{independence}: in \textsc{const} the observed
        ``both silenced'' rate is $0.1763$ against $0.1775$ predicted by eq.~\eqref{eq:G2}, but against
        $\Hbar^2=0.0625$ if losses were not inherited --- a miss by $z=34.0$;
  \item \textsc{early}'s events scored with \textsc{const}'s prediction: in clades of $5$--$16$ cells,
        $8.8$ visible events per tree observed against $29.9$ under constant timing ($8.5$ under the
        right timing) --- a miss by $z=45.7$.
\end{itemize}
G2 and G3 can see both inheritance and timing.

\paragraph{Disclosed.} The per-tape cells rest on a few rare clone-wide events (about $6$--$19$ stem
events per tape over 200 trees), so their $t$-statistics sit near the edge of calibration: their spread
was $1.28$ (G1 \textsc{const}) and $1.35$ (G5) at the first seed against $\approx1.0\pm0.13$, and $0.91$
and $1.18$ at the passing seed. No bias was seen at either seed.

\paragraph{What the fixtures look like} (descriptive, not science; $\Hbar=0.25$, $n=210$, the
validation's $\capt\times\turn$ mix). Heritable share $75\%$ \textsc{const}, $86\%$ \textsc{early}, $31\%$
\textsc{late}; visible events per tape per tree $17.6$, $9.0$, $41.8$; spread of the realised $\Hbar$
across trees $0.036$, $0.059$, $0.008$. Early silencing makes the realised fraction about $7\times$
noisier than late, as eq.~\eqref{eq:Hvar} predicts.

\paragraph{Cost.} $28$--$45$~s on one core for the whole validation, $0.46$~GB. One $\Exp(1)$ per
(branch, tape) is $(2n-1)k\approx1.3\times10^4$ draws per tree at $n=210$.

\subsection{The rung-2 library}\label{sec:drop2lib}

Built on the rung-1 library's $30$ tree cells and $5$ panel replicates, with $21$ settings:
$\Hbar$ = $0$, $0.01$, $0.02$, $0.05$, $0.10$, $0.20$, $0.33$, crossed with three timings --- constant;
\emph{early}, $\zeta_0=2$ on $[0,0.4)$ and $\tfrac13$ after; and \emph{late}, $\zeta_0=\tfrac13$ on
$[0,0.6)$ and $2$ after. \emph{Early} puts
$80\%$ of the hazard in the first $40\%$ of the experiment, \emph{late} $80\%$ in the last $40\%$. Loss
draws are per (tree, replicate) and shared by every $\LamT$, panel, mask and setting. Because sets are
nested, \emph{one} sparse record per tree and shape --- every (branch, tape, $E$) with $E\le\Gamma_e$ at
$\Hbar=0.33$ --- rebuilds $H$ at any lower $\Hbar$. A dataset is composed at read time from the rung-1
file at the same (replicate, tree): $Y=\neg X\wedge\neg H\wedge(D\ge1)$. \emph{Result:} $134$~MB, zero
assertion violations; realised minus target $\Hbar$ averages $+0.0002$ to $+0.0010$ at every $n$.

\begin{center}\small
\emph{Reading rule:} each entry is the expected heritable share, eq.~\eqref{eq:hershare} --- the
percentage of silenced entries that sit in a clade of at least two cells --- at $\Hbar=0.10$,
$\turn=0.3$, for constant / early / late timing. $0\%$ means all silencing is in singletons and looks
like technical dropout; higher means more of it is visible to a relatedness-based statistic.\\[3pt]
\begin{tabular}{@{}lccc@{}}
\toprule
cells $n$ & $\capt=0.002$ & $\capt=0.1$ & $\capt=0.25$ \\
\midrule
25      & 50 / 80 / 20\% & 75 / 91 / 53\% & 81 / 94 / 63\% \\
210     & 59 / 86 / 27\% & 81 / 94 / 63\% & 87 / 96 / 74\% \\
1{,}976 & 66 / 89 / 34\% & 86 / 95 / 71\% & 90 / 97 / 80\% \\
\bottomrule
\end{tabular}
\end{center}
\noindent The heritable share rises with capture, with clade size, and with earlier timing, from $20\%$
($n=25$, $\capt=0.002$, late) to $97\%$ ($n=1{,}976$, $\capt=0.25$, early). Capture matters because low
$\capt$ means long terminal branches (Fig.~\ref{fig:S3ab}), and a long terminal branch collects
hazard that only one cell can show.

\paragraph{A correction to the design.} The rung-2 proposal predicted that long terminal branches at low
capture would push \emph{most} losses into singletons. Under constant timing the singleton share at
$\capt=0.002$ is $34$--$50\%$, a minority; it is a majority only under late timing ($66$--$80\%$). The
direction was right --- the heritable share falls with $\capt$ in every row --- but the magnitude was
overstated. No recorded number moves.

\paragraph{One small bias, untested.} The \emph{realised} heritable share, a per-tree ratio averaged over
trees, sits $0.5$--$2.5$ points below the expected one in all nine $n=210$ rows. Eq.~\eqref{eq:hershare}
is a ratio of expectations, and its numerator and denominator co-vary (one big event raises both).
Validation checks the expectations (G3), not this ratio.

\subsection{What the two rungs do not represent}\label{sec:dropnot}

Rung~2, like rung~1, generates a \emph{model}, and simulated data are cleaner than Park's in known ways:
\begin{itemize}
  \item \textbf{Complete loss only.} The rebuilt event detector puts Mouse3's partial-loss fraction at
        $29.9\%$; in rung~2 every loss is complete, so a completeness estimate reads $1$ by construction.
        Partial loss would be a later switch.
  \item \textbf{No \emph{trans} limb.} A silenced integration's pegRNA symbol keeps being written on other
        tapes, so rung~2 cannot validate the symbol-depletion test.
  \item \textbf{Silencing per unit time.} If silencing actually happens at replication, the hazard along
        a reconstructed branch follows the lineage's division rate, which is not constant. A shaped
        $\zeta_0$ approximates that but does not represent it.
  \item \textbf{The same propensity on every tape.} Park cannot yet say otherwise: a tape's low recovery
        (park-compatibility Fig.~3b) is read as its integration's expression level, which fits both
        ``prone to silencing'' and ``already silenced in the shared founder line'', and a naive per-tape
        event count is biased because a tape with high baseline missingness gives a loss less contrast.
        If propensity does vary, simulated losses are spread too evenly across tapes.
  \item \textbf{One tree, one clone, no shared ancestry.} A multi-clone arm descends from one engineered
        line, and Park's clone-wide loss layer is read as pre-bottleneck silencing. This build cannot
        make a tape lost in several clones for a shared reason.
  \item \textbf{QC selection.} Silencing lowers $R_c$, so $R_{\min}$ preferentially removes heavily
        silenced cells. G1--G5 are scored before the filters, as in rung~1.
\end{itemize}
Nor is anything here calibrated to Park. Park's event catalogue is detection-limited (no events in
clones under 20 cells, counts that depend on the threshold, Pre-TX possibly inflated by censoring), and
its silenced shares --- $2.46$--$20.07\%$ of missing entries at the budget criterion, plus a clone-wide
layer of $7$--$33\%$ of missing at a $10$-nat floor --- serve only to orient the grid: a share of missing
is at most a share of entries, so a grid reaching $\Hbar=0.33$ covers them. Calibration means pushing
rung-2 data through the same caller and matching its output, which is the next step: the negative
control as a dose--response, rung-1 and rung-2 data through \texttt{73}, \texttt{74}, \texttt{77} and the
event caller across $\Hbar\times\zeta_0$.
